\begin{frame}{Ontologies, meaning, and the open world assumption}

  \begin{itemize}
    \item A \textbf{vocabulary model} is a community agreement: \emph{what the sets contain}.
    \begin{itemize}
      \item \texttt{Terminal}, \texttt{ConnectivityNode}, \texttt{AnalogValue}, \texttt{PhaseCode} — reusable across use cases, across organizations.
    \end{itemize}
    \pause
    \item The \textbf{Open World Assumption (OWA)}:
    \begin{itemize}
      \item what data does \emph{not} say is \emph{not false} — someone else may say more
      \item different data can say \emph{more} about the same world without contradicting each other
      \item this is what makes \emph{integration} of different standards possible at all
    \end{itemize}
    \pause
    \item The \textbf{Closed World Assumption (CWA)} allows us to validate data against rules:
    \begin{itemize}
      \item a profile / schema says: \emph{this is what the data is supposed to look like}
      \item this is what makes \emph{validation} and \emph{interoperability of implementations} possible
    \end{itemize}
  \end{itemize}

\end{frame}
